\section{Alphabet and grammar}
\label{sec:prop-syntax-semantics-01}

\begin{definition}[Propositional alphabet]
The \emph{propositional alphabet} consists of:
\begin{enumerate}
  \item propositional variables
        \(\Prop\defeq\set{p_n:n\in\N}\);
  \item the nullary connective \(\bottom\);
  \item the unary connective \(\neg\);
  \item the binary connectives \(\land,\lor,\to\); and
  \item the punctuation symbols \((\) and \()\).
\end{enumerate}
All listed symbols are distinct.  A propositional variable is an
object-language symbol.  The letters \(\varphi,\psi,\chi\) used below are
metavariables ranging over formulas.  For readable examples, \(P,Q,R,S\)
denote specified members of \(\Prop\); the indexed names \(p_n\) are used when
the identity of a variable matters.
\end{definition}

\begin{definition}[Propositional formula]
The set \(\Fm_{\CPL}\) of \emph{propositional formulas} is the least set of
strings satisfying the following formation clauses:
\begin{enumerate}
  \item \(p_n\in\Fm_{\CPL}\) for every \(n\in\N\);
  \item \(\bottom\in\Fm_{\CPL}\);
  \item if \(\varphi\in\Fm_{\CPL}\), then
        \(\neg\varphi\in\Fm_{\CPL}\); and
  \item if \(\varphi,\psi\in\Fm_{\CPL}\) and
        \(\circ\in\set{\land,\lor,\to}\), then
        \((\varphi\circ\psi)\in\Fm_{\CPL}\).
\end{enumerate}
No other string is a propositional formula.
\end{definition}

The displayed grammar is
\[
  \varphi ::= p_n\mid\bottom\mid\neg\varphi
  \mid(\varphi\land\varphi)\mid(\varphi\lor\varphi)
  \mid(\varphi\to\varphi).
\]
Parentheses required by the grammar are part of the object language.  In prose
and displays, outermost parentheses may be omitted, and the precedence order
is \(\neg\), \(\land\), \(\lor\), \(\to\).  A chain of repeated
\(\land\), repeated \(\lor\), or repeated \(\to\) associates to the right.
These are typographical abbreviations, not additional formation clauses.

\begin{definition}[Defined connectives]
The symbols \(\topv\) and \(\leftrightarrow\) are metalinguistic
abbreviations:
\[
  \topv\defeq\neg\bottom,
  \qquad
  (\varphi\leftrightarrow\psi)
  \defeq((\varphi\to\psi)\land(\psi\to\varphi)).
\]
They do not add constructors to \(\Fm_{\CPL}\).
\end{definition}

For an ordered finite sequence \(\varphi_1,\ldots,\varphi_n\), define its
finite conjunction and disjunction as right-associated formulas.  When
\(n>0\), put \(C_n=\varphi_n\), \(D_n=\varphi_n\), and recursively put
\[
 C_i=(\varphi_i\land C_{i+1}),
 \qquad
 D_i=(\varphi_i\lor D_{i+1})
 \quad(1\leq i<n);
\]
then
\[
 \bigwedge_{i=1}^{n}\varphi_i\defeq C_1,
 \qquad
 \bigvee_{i=1}^{n}\varphi_i\defeq D_1.
\]
For \(n=0\), define the finite conjunction as \(\topv\) and the finite
disjunction as \(\bottom\).  An indexed finite conjunction or disjunction is
used only after its index order has been stated.  This convention applies to
both propositional and first-order formulas and gives every later finite
connective display one formula-string value.

\begin{lemma}[Parenthesis counting]
\label{lem:paren-count}
Every member of \(\Fm_{\CPL}\) contains equally many left and right
parenthesis occurrences.  Moreover, for every formula and every
binary-connective occurrence in it, the number of left parenthesis
occurrences strictly preceding that occurrence exceeds the number of right
parenthesis occurrences strictly preceding it.
\end{lemma}
\begin{proof}
Both claims proceed by structural induction on the formula.  A propositional
variable and \(\bottom\) contain no parentheses and no binary-connective
occurrences.  The negation clause prefixes one token and adds no
parentheses, so both claims pass through it unchanged.  For
\((\varphi\circ\psi)\), balance holds because the clause adds one left and
one right parenthesis to balanced parts.  A binary-connective occurrence in
\((\varphi\circ\psi)\) is the displayed occurrence of \(\circ\), an
occurrence in \(\varphi\), or an occurrence in \(\psi\).  The tokens
preceding the displayed \(\circ\) are the initial left parenthesis followed
by the balanced string \(\varphi\), so the excess of left over right
parentheses there is one.  For an occurrence in \(\varphi\), the initial
left parenthesis raises the positive excess supplied by the induction
hypothesis by one.  For an occurrence in \(\psi\), the preceding tokens are
the initial left parenthesis, the balanced string \(\varphi\), the
connective, and a prefix of \(\psi\), so the excess is one more than the
positive excess supplied by the induction hypothesis within \(\psi\).
\end{proof}

\begin{theorem}[Unique readability]
\label{thm:prop-unique-readability}
Every member of \(\Fm_{\CPL}\) satisfies exactly one formation clause.  In
the negation case its immediate subformula is unique.  In a binary case its
main connective and its two immediate subformulas are unique.
\end{theorem}
\begin{proof}
A propositional variable begins with a variable token, \(\bottom\) is a
distinct one-token string, a negation begins with \(\neg\), and every binary
formula begins with a left parenthesis.  Hence the four kinds of formation
clause have disjoint ranges.

For a binary formula \((\varphi\circ\psi)\), scan from immediately after
its first left parenthesis while maintaining the number of left minus right
parenthesis occurrences already scanned.  By the balance clause of
\cref{lem:paren-count}, this number is zero at the displayed occurrence of
\(\circ\), because the tokens scanned before it form exactly the balanced
string \(\varphi\).  By the excess clause of \cref{lem:paren-count},
applied to \(\varphi\) and to \(\psi\), every binary connective inside
either immediate subformula occurs while the number is positive.  Hence the
count-zero binary-connective position is unique.  The first formula ends
exactly before that unique position, and the second begins immediately after
it and ends immediately before the final right parenthesis.  Thus the split, the
main connective, and the two immediate subformulas are unique.  Removing an
initial \(\neg\) uniquely determines the immediate subformula in the unary
case.  Therefore every formula has one and only one outermost construction.
\end{proof}

\begin{definition}[Complexity and subformula]
The \emph{complexity} \(c(\varphi)\in\N\) is defined by
\[
\begin{aligned}
 c(p_n)&=0,&c(\bottom)&=0,\\
 c(\neg\varphi)&=c(\varphi)+1,
 &c((\varphi\circ\psi))&=c(\varphi)+c(\psi)+1.
\end{aligned}
\]
The set \(\Sub(\varphi)\) of \emph{subformulas} is defined by
\[
\begin{aligned}
 \Sub(p_n)&=\set{p_n},&\Sub(\bottom)&=\set{\bottom},\\
 \Sub(\neg\varphi)&=\set{\neg\varphi}\cup\Sub(\varphi),\\
 \Sub(\varphi\circ\psi)&=
 \set{\varphi\circ\psi}\cup\Sub(\varphi)\cup\Sub(\psi).
\end{aligned}
\]
Here and below \(\circ\in\set{\land,\lor,\to}\).
\end{definition}

\begin{hbdefinition}{def:p02-w02-p02-044-syntax-tree}{Syntax tree of a propositional formula}
The \emph{syntax tree} of a formula is the finite rooted tree, with ordered
children and labelled nodes, given by recursion on the formation clauses.
The syntax tree of \(p_n\) is one node labelled \(p_n\), and the syntax tree
of \(\bottom\) is one node labelled \(\bottom\).  The syntax tree of
\(\neg\varphi\) has a root labelled \(\neg\) whose only child is the root of
the syntax tree of \(\varphi\).  The syntax tree of \((\varphi\circ\psi)\) has
a root labelled \(\circ\) whose left child is the root of the syntax tree of
\(\varphi\) and whose right child is the root of the syntax tree of \(\psi\).
The recursion assigns exactly one tree to each formula because, by
\cref{thm:prop-unique-readability}, each formula is produced by exactly one
formation clause, and the main connective and immediate subformulas used by
that clause are unique.  The leaves are the occurrences of propositional
variables and of \(\bottom\); the internal nodes are the connective
occurrences.
\end{hbdefinition}

The display \(\neg P\lor Q\land R\to S\to\bottom\) abbreviates, under the
precedence order and the right association of \(\to\) stated after the
grammar, the formula
\[
  ((\neg P\lor(Q\land R))\to(S\to\bottom)).
\]
Scanning from immediately after its first left parenthesis, as in the proof
of \cref{thm:prop-unique-readability}, the number of left minus right
parentheses already scanned is \(1\) at \(\lor\), \(2\) at \(\land\), \(0\)
at the first \(\to\), and \(1\) at the second \(\to\).  The first \(\to\) is
therefore the main connective, with immediate subformulas
\((\neg P\lor(Q\land R))\) and \((S\to\bottom)\).  Splitting each of these in
the same way gives the syntax tree in \cref{fig:p02-w02-p02-044-syntax-tree}.
Its ten nodes are the roots of the syntax trees of the ten members of the
subformula set: the whole formula, \((\neg P\lor(Q\land R))\), \(\neg P\),
\(P\), \((Q\land R)\), \(Q\), \(R\), \((S\to\bottom)\), \(S\), and
\(\bottom\).  The complexity of the formula is \(5\), one for each of its
five connective occurrences.  The tree has \(5\) leaves and \(4\)
binary-connective nodes; the negation node has one child and does not change
that difference.

\begin{figure}[tp]
\centering
\begin{tikzpicture}[level distance=9mm,
 every node/.style={inner sep=1.5pt},
 level 1/.style={sibling distance=36mm},
 level 2/.style={sibling distance=16mm},
 level 3/.style={sibling distance=9mm}]
\node {\(\to\)}
 child {node {\(\lor\)}
   child {node {\(\neg\)} child {node {\(P\)}}}
   child {node {\(\land\)}
     child {node {\(Q\)}}
     child {node {\(R\)}}}}
 child {node {\(\to\)}
   child {node {\(S\)}}
   child {node {\(\bottom\)}}};
\end{tikzpicture}
\caption{The syntax tree of \(((\neg P\lor(Q\land R))\to(S\to\bottom))\).
Internal nodes are connective occurrences, leaves are occurrences of
propositional variables and of \(\bottom\), and the subtree below each node is
the syntax tree of one subformula.  The root is the unique count-zero
connective found in the proof of \cref{thm:prop-unique-readability}.}
\label{fig:p02-w02-p02-044-syntax-tree}
\end{figure}

